\documentclass[UTF8,10.9pt]{ctexart}
\usepackage{lesson-style}
% Use after lesson-style.sty: source-inspired blue vector nodes and arrows.
\tikzset{
 lessonnode/.style={circle,draw=LessonGraph,very thin,fill=LessonPaleBlue,minimum size=8.5mm,inner sep=1pt,font=\small\bfseries,text=LessonBlue},
 lessonedge/.style={-{Stealth[length=2.6mm]},line width=.8pt,draw=LessonGraph},
 lessonedgemuted/.style={-{Stealth[length=2.6mm]},line width=.75pt,draw=LessonBlueStripe},
 lessonnote/.style={font=\footnotesize,text=LessonTeal,align=left},
 lessonemphasis/.style={font=\footnotesize,text=LessonOrange,align=left},
 lessonbox/.style={draw=LessonBlueStripe,fill=LessonPaleBlue,rounded corners=1pt,inner sep=5pt}
}

\begin{document}
\LessonTitle{拓扑排序}{有向图的先后关系、队列实现与课堂练习}
\section{一、知识点讲解}
\subsection{1. 什么是拓扑排序？}
很多事情都有先后顺序：先学习基础课程，才能学习提高课程。把每件事情看作顶点，把“$u$ 必须先于 $v$”画成有向边 $u\to v$。拓扑排序要求每条有向边的起点都排在终点前面。
\begin{lessonhint}
\textbf{拓扑排序必须满足的规则}\par
如果存在 $u\to v$，则输出序列中 $u$ 必须出现在 $v$ 前面。合法序列不一定唯一。
\end{lessonhint}
\begin{center}
\begin{tikzpicture}[node distance=15mm and 25mm]
\node[lessonnode] (a) {1};
\node[lessonnode,right=of a,yshift=10mm] (b) {2};
\node[lessonnode,right=of a,yshift=-10mm] (c) {3};
\node[lessonnode,right=of b,yshift=-10mm] (d) {4};
\draw[lessonedge] (a)--(b);
\draw[lessonedge] (a)--(c);
\draw[lessonedge] (b)--(d);
\draw[lessonedge] (c)--(d);
\end{tikzpicture}
\end{center}
\noindent 合法顺序包括 $1,2,3,4$ 和 $1,3,2,4$。
\subsection{2. 入度、出度与 DAG}
\begin{center}
\begin{tabularx}{.96\linewidth}{@{}lX@{}}\toprule
概念 & 含义\\\midrule
入度 $indeg[v]$ & 指向顶点 $v$ 的有向边数量\\
出度 & 从顶点出发的有向边数量\\
入度为 0 & 没有尚未处理的前置要求\\
DAG & 不包含有向环的有向图\\\bottomrule
\end{tabularx}
\end{center}
\subsection{3. Kahn 算法：队列实现}
\begin{enumerate}
\item 统计每个顶点的入度，并建立邻接表。
\item 把所有入度为 0 的顶点放入队列。
\item 取出队首顶点 $u$，加入答案。
\item 对每条边 $u\to v$，执行 $indeg[v]--$；若变为 0，则让 $v$ 入队。
\item 队列为空后，若处理数量少于 $n$，则图中存在有向环。
\end{enumerate}
\LessonLabel{核心代码：队列与入度更新}
\begin{lstlisting}[style=lessonfragment]
while (!q.empty()) {
    int u = q.front();
    q.pop();
    ans.push_back(u);
    for (int v : g[u]) {
        indeg[v]--;
        if (indeg[v] == 0) q.push(v);
    }
}
\end{lstlisting}
\section{二、课堂练习}
\subsection{1. 示例：任务先后顺序}
\noindent 题目 ID：示例题\qquad 时间限制：1000 ms\qquad 空间限制：256 MB
\LessonLabel{题目描述}
给定 $n$ 个任务和 $m$ 条依赖关系 $u\to v$，输出任意一种合法任务执行顺序。保证存在合法答案。
\LessonLabel{输入格式}
第一行两个整数 $n,m$；接下来 $m$ 行，每行两个整数 $u,v$。
\LessonLabel{输出格式}
输出一个合法拓扑序。
\LessonLabel{样例 1}
\begin{center}\begin{tabularx}{.9\linewidth}{XX}\toprule 输入 & 输出\\\midrule
\texttt{4 4} & \texttt{1 2 3 4}\\
\texttt{1 2} & \\
\texttt{1 3} & \\
\texttt{2 4} & \\
\texttt{3 4} & \\\bottomrule\end{tabularx}\end{center}
\LessonLabel{题意概括}
任务是顶点，依赖关系是有向边；输出所有顶点的一种拓扑序。
\LessonLabel{思路分析}
第一步：建图并统计入度。第二步：把入度为 0 的任务入队。第三步：反复取出任务，并减少后继任务的剩余入度。
\LessonLabel{样例推导分析过程}
\begin{center}\begin{tabularx}{.95\linewidth}{lXX}\toprule 步骤 & 当前处理 & 处理后的队列\\\midrule
初始 & — & [1]\\
1 & 1 & [2,3]\\
2 & 2 & [3]\\
3 & 3 & [4]\\
4 & 4 & 空\\\bottomrule\end{tabularx}\end{center}
\LessonLabel{复杂度分析}
时间复杂度 $O(n+m)$；空间复杂度 $O(n+m)$。
\LessonLabel{AC 代码（示例实现，未连接在线评测）}
\begin{lstlisting}
#include <bits/stdc++.h>
using namespace std;
int main() {
    ios::sync_with_stdio(false);
    cin.tie(nullptr);
    int n, m;
    cin >> n >> m;
    vector<vector<int>> g(n + 1);
    vector<int> indeg(n + 1, 0);
    for (int i = 0; i < m; i++) {
        int u, v;
        cin >> u >> v;
        g[u].push_back(v);
        indeg[v]++;
    }
    queue<int> q;
    for (int i = 1; i <= n; i++)
        if (indeg[i] == 0) q.push(i);
    while (!q.empty()) {
        int u = q.front(); q.pop();
        cout << u << ' ';
        for (int v : g[u]) {
            indeg[v]--;
            if (indeg[v] == 0) q.push(v);
        }
    }
    cout << '\n';
    return 0;
}
\end{lstlisting}
\end{document}
